\section{Introduction}
Brain tumors pose a significant challenge in clinical practice due to their diverse appearances, growth patterns, and locations within the brain \cite{1}. Among the various intracranial growths encountered, arteriovenous malformations (AVMs), pituitary tumors, meningiomas, schwannomas, and metastatic lesions account for a large percentage of diagnosed cases. Although Magnetic Resonance Imaging (MRI) is the established standard for evaluating brain tumors, providing high-resolution details of soft tissues and anatomy \cite{2}, distinguishing between different tumor types can be difficult. The precise classification and segmentation of these tumors remains challenging due to overlapping features seen on imaging and the variability within individual tumors \cite{3,4}. The necessity of precise diagnostic tools that can accurately differentiate between these tumor types underscores the importance of innovations in medical imaging technology \cite{5,6}. While histopathology and molecular analysis remain essential for definitive diagnosis, prognostication, and treatment planning in neuro-oncology, these methods are time-consuming, require invasive procedures, and are not always practical in urgent clinical situations. These limitations have led to a growing interest in automated, high-precision diagnostic tools that can deliver reliable results with minimal human intervention. In this work, we aim to enhance non-invasive imaging analysis prior to histopathology.

Each type of brain tumor presents unique challenges for both imaging and biological interpretation. AVMs, being congenital vascular abnormalities, have complex, tangled vessel structures that make them difficult to outline precisely. Pituitary tumors, though often small, exhibit a wide range of functional profiles that impact hormone levels and thus complicate diagnosis. Meningiomas and schwannomas can appear similar on MRI scans despite originating from different cell types. The diverse imaging characteristics of metastases, which depend on their primary source, add another layer of complexity. Accurately classifying and delineating these various tumor types is crucial for effective treatment planning, guiding surgical procedures, and predicting patient outcomes.

Convolutional neural networks (CNNs) have shown strong performance in both classification and segmentation tasks, including lung cancer classification \cite{7} and brain tumor analysis \cite{8,9,10}. Notably, models like ResBRNet and EfficientNet-B0 have achieved impressive results when trained on large-scale datasets \cite{11}. Similarly, architectures such as U-Net \cite{12} and NN U-Net \cite{13} have become the standard for tumor segmentation tasks. DeepLabV3+ and Efficient-Enhanced U-Net also achieved impressive results when trained on BraTS dataset \cite{14,15}. However, CNNs inherently struggle with spatial variations, especially rotations and translations. This requires extensive data augmentation, which can lead to inconsistencies in how features are extracted from MRIs acquired in different orientations. Consequently, this increases computational demands and reduces the interpretability of the models. Although transfer learning has improved performance in some instances \cite{16}, pre-trained models still lack an inherent understanding of geometric invariance.

To overcome these limitations, geometric deep learning (GDL) has emerged as a promising alternative. Models such as Group Equivariant CNNs (G-CNNs) \cite{17} and Harmonic Networks \cite{18} are designed to respect rotational symmetry, making them more robust to changes in orientation. However, many of these models primarily focus on rotational equivariance and often overlook translational invariance, which is vital for accurately locating tumors. Previous research has typically concentrated on either classification or segmentation, rarely addressing both tasks within a unified model or encompassing the full spectrum of brain tumor diversity.

GDL provides a biologically inspired approach to incorporating spatial knowledge and structural symmetry into neural networks \cite{19,20,21}. It has gained increasing traction across medical imaging domains beyond brain tumor classification. For instance, Bronstein et al. demonstrated the GDL in analyzing cortical surfaces and mapping functional regions using non-Euclidean mesh representations \cite{22}. In diffusion MRI, Ewert et al. introduced a GDL method for continuous signal reconstruction \cite{23}. A deeper representation of GDL is also employed. Zhang et al. applied local-to-global graph neural networks to classify Autism Spectrum Disorder (ASD) and Alzheimer's disease (AD) by analyzing disease-related local brain regions and biomarkers \cite{24}. Additionally, Kerkela et al. modeled spherical neural networks to improve the estimation of microstructural data from brain MRI data \cite{25}. This research improves prediction accuracy compared to models that apply less or no rotational variance. This emphasizes the effectiveness of a model in medical images for models that use rotational features, such as GDL methods. Zhu et al. proposed a unified GDL model to predict task-based fMRI activity from anatomical features in sMRI, highlighting cross-modal learning \cite{26}. In CT imaging, Koetzier et al. presented a graph-based deep reconstruction method to recover high-fidelity CT volumes from undersampled data \cite{27}. Winkels and Cohen proved CNN model with roto-translation group convolutions performs similarly to a regular CNN with ten times less data \cite{28}. This research trained a roto-translation group 3D CNN on a pulmonary nodule detection dataset and compared it with a regular CNN model. The regular CNN model received data augmentation, yet still needed ten times more data to perform at an acceptable level. In digital pathology, Li et al. use GDL to conduct spatial analysis using histopathology images for single-cell analysis \cite{29}. Biologically, analyzing histopathology images at the single-cell level offers a significant advantage over traditional patch-based methods by providing detailed cellular information. This research focuses on exploring more into the microenvironment of a tumor. In their experiments, their proposed model has consistently eclipsed existing patch-based and cell graph-based methods across two independent datasets.

This study aims to address two fundamental tasks in neuro-oncological imaging, which are classifying various tumor types and performing semantic segmentation. For classification, we use a private MRI dataset containing five distinct tumor categories: AVM, meningioma, pituitary adenoma, metastatic tumors, and schwannoma. Arteriovenous malformations are vascular malformations with an incidence of approximately 1 per 100,000 per year and carry a risk of serious hemorrhage, where their detection via MRI is critical for preventing stroke and guiding surgical planning \cite{30,31}. Meningiomas account for roughly 30-37\% of all primary brain tumors, are frequently benign, and often produce symptoms through mass effect, with MRI guiding both diagnosis and surgical strategy \cite{32}. Pituitary adenomas represent 10–25\% of intracranial neoplasms, with microadenomas being particularly common (prevalence up to $\sim$17\%), and can cause visual or hormonal symptoms that MRI is uniquely suited to identify \cite{33}. Brain metastases, occurring in an estimated 10–40\% of cancer patients, are the most frequent intracranial tumors, significantly impacting prognosis and treatment decisions \cite{34}. Vestibular schwannomas (also called acoustic neuromas) are benign tumors arising from Schwann cells of the vestibular portion of the eighth cranial nerve. They account for approximately 8–10\% of all intracranial tumors and up to 80\% of tumors in the cerebellopontine angle. Patients typically present with hearing loss, tinnitus, and balance disturbances \cite{35}. Collectively, these tumor types underscore the need for an imaging tool capable of both classification and localization across heterogeneous tumor morphologies in real-world neuro-oncology practice. For segmentation, we perform pixel-level labeling to precisely outline tumor regions. This is essential for estimating tumor volume and tracking its boundaries. Our goal is to provide radiologists with clinically valuable insights to assist them in diagnosis and treatment planning.

To achieve these goals, we introduce an upgraded Mod-SE(2) pipeline:
1) A \textbf{2.5D Spatial Context Model} that processes $n-1, n, n+1$ slices to encode structural continuity across the $z$-axis.
2) An \textbf{SE(2)-Equivariant CNN (SE2-CNNET)} that structurally guarantees roto-translational equivariance, ensuring robust feature extraction across varying orientations.
3) An advanced \textbf{Edge-Boundary Loss} function that rigorously enforces accurate boundary delineation by penalizing boundary misclassifications utilizing morphological operations.

Our primary contributions include demonstrating improved performance compared to conventional models across public and private datasets, and proving that incorporating 2.5D spatial awareness with geometric priors results in highly accurate and clinically interpretable segmentation masks while maintaining computational efficiency.
